\documentclass[11pt,letterpaper]{article}
\usepackage[T1]{fontenc}
\usepackage[utf8]{inputenc}
\usepackage{lmodern,microtype,amsmath,booktabs,tabularx,longtable,graphicx}
\usepackage[margin=0.8in,headheight=15pt]{geometry}
\usepackage{xcolor,fancyhdr,enumitem,xurl}
\usepackage[hidelinks]{hyperref}
\definecolor{Gold}{HTML}{987100}
\setlength{\parindent}{0pt}\setlength{\parskip}{6pt}
\setlist{nosep,leftmargin=*}\emergencystretch=3em
\pagestyle{fancy}\fancyhf{}\fancyhead[L]{\small Follow My Voice}\fancyfoot[R]{\thepage}
\newcommand{\qid}[1]{\par\smallskip\textbf{\texttt{#1}}\par}
\begin{document}
\setcounter{section}{0}
\section{Introduction}
Finding an unnamed target in a surrounding mixed-reality scene requires a participant to coordinate gaze, head and chair rotation, and controller selection. Spoken assistance can reveal the target through a sequence of directions while leaving the participant responsible for acting on those directions. The voice delivering this assistance may shape both performance and the experience of being guided.

Follow My Voice will compare a participant-selected catalog voice with a synthetic version of the participant's own voice. Both conditions will use the same gaze-contingent guidance policy. Participants will remain seated among eight surrounding object planes and select objects with the right-controller ray and trigger. The agent will know the target; the participant will discover it through guidance rather than a prior color-and-shape description.

The primary research question is whether voice condition changes active search-completion time after accounting for the target's horizontal angle from the trial starting reference. Secondary questions concern workload and motivation. Exploratory questions concern perceived agent presence, intelligence, humanlike qualities, voice experience and reliance. The within-participant comparison will use two counterbalanced blocks of 40 measured trials. No directional performance benefit is assumed.

The voice contrast concerns two personalized experiences: a preferred catalog voice and a self-similar voice. It will not isolate the effect of gaze contingency because both conditions receive it. Reports that a voice resembles inner thought or feels attentive will be treated as participant interpretations rather than evidence of a particular cognitive mechanism or agent sentience.
\end{document}
